\chapter{Remaining FYP Plan}\label{ch:remaining}

\section{Remaining Scope}
The remaining scope completes the satellite-to-GIS chain on top of the platform built in Iteration 1:
the selected crop-type pipeline in production use, a dedicated field-boundary model with area
validation, yield estimation, wider evaluation across areas, and a hardened, deployed system.

\section{Remaining Modules and Features}
\begin{table}[htbp]
\centering\small
\caption{Remaining modules and features}
\label{tab:remaining}
\begin{tabularx}{\textwidth}{L{3.6cm} Y L{1.8cm}}
\toprule
\textbf{Module} & \textbf{Feature} & \textbf{Iteration}\\
\midrule
M1 Data preparation & Monthly-composite ingestion in the pipeline; exports for remaining AOIs; optional cloud mask & 2, 3\\
M2 Crop mapping & LightGBM crop-type pipeline as a registry model; CDL confidence rule; cross-area test; field-level classification & 2, 3\\
M3 Boundaries and area & Field-boundary model; polygon simplification; IoU and area error; acreage against NASS & 3\\
M4 Yield & NASS yield preparation; yield model; county and field summaries & 3, 4\\
M5 Processing pipeline & Multi-image upload; batch progress; mosaics & 2\\
M8 Web GIS portal & All-areas view; compare runs; yield layers and dashboard; exports (GeoJSON, CSV) & 2--4\\
Hardening & Required sign-in, signed tile access, monitoring, load test, provider basemap & 4\\
\bottomrule
\end{tabularx}
\end{table}

\section{Remaining Iterations}
\begin{itemize}
  \item \textbf{Iteration 2 (mid-Oct -- Dec 2026, FYP-1 end):} crop-type pipeline integration, CDL
        confidence rule, AOI2 interval, cross-area test, multi-image upload and all-areas view,
        serving-cost measurement.
  \item \textbf{Iteration 3 (Feb -- Mar 2027, FYP-2):} field-boundary model and area validation; yield
        data and first yield model; more AOIs; field-level crop classification experiment.
  \item \textbf{Iteration 4 (Apr -- May 2027, FYP-2):} yield layers and dashboard, end-to-end evaluation
        against the success criteria, hardening, deployment, final report and demonstration.
\end{itemize}

\section{Implementation Plan}
New models are added through the registry (task and manifest), so the API, schema and interface
remain stable. The LightGBM pipeline is packaged with its saved model, feature definition and
metrics file; the ingest step gains a monthly-composite path. The boundary model writes a raster
or polygons that the existing vector-tile builder publishes as PMTiles. The yield model reads the
crop-type layer and satellite features and writes a continuous raster plus county summaries. The
portal adds layer groups and dashboard panels driven by the catalogue.

\section{Testing and Validation Plan}
\begin{itemize}
  \item Unit tests for each new model task (validation, outputs, metrics) and the extended pipeline.
  \item Integration and end-to-end tests repeated for every release; benchmark re-run with the new
        model and with several areas and layers visible.
  \item Crop type: all-pixel and CDL-confidence scores, cross-area test, per-class scores, intervals.
  \item Boundaries: IoU/Dice against reference boundaries and area error; acreage against USDA NASS.
  \item Yield: $R^2$, RMSE and MAE on held-out years or counties.
  \item Usability study (P3): five to eight participants, including the supervisors and, if
        possible, an agricultural user, complete set tasks (find a crop, compare two layers, upload
        an image); task success, time and a System Usability Scale questionnaire are recorded.
  \item Resource measurement repeated for every new model (time, peak RAM, storage), as in
        Table~\ref{tab:resources}.
\end{itemize}

\section{Deployment Plan}
The same Docker Compose stack will be deployed to a single cloud virtual machine for the final
demonstration, with a provider basemap key, required sign-in, HTTPS through the edge server and
persistent volumes. The architecture allows managed PostgreSQL, S3 storage, a CDN in front of the
tile path and additional worker and tile replicas if more users or areas are added.

\section{Tentative Timeline}
\begin{table}[htbp]
\centering\small
\caption{Tentative timeline}
\label{tab:timeline}
\begin{tabularx}{\textwidth}{L{3.2cm} Y}
\toprule
\textbf{Period} & \textbf{Planned work}\\
\midrule
Oct 2026 & Mid-term demonstration; start crop-type pipeline integration\\
Nov 2026 & Pipeline integrated; CDL confidence rule; cross-area test; multi-image upload\\
Dec 2026 & Serving-cost measurement; FYP-1 final report and demonstration\\
Feb 2027 & Boundary model; remaining AOI exports; yield data preparation\\
Mar 2027 & Boundary and area validation; first yield model\\
Apr 2027 & Yield in portal; end-to-end evaluation; hardening; deployment\\
May 2027 & Final evaluation, report and demonstration\\
\bottomrule
\end{tabularx}
\end{table}

\section{Expected Final Deliverables}
\begin{itemize}
  \item Deployed GeoCrop platform with crop-type, field-boundary and yield layers, dashboard and model
        catalogue.
  \item Source code repositories with tests, documentation and reproducible data preparation.
  \item Trained models with model cards and stored metrics.
  \item Evaluation report covering crop type, boundaries, area and yield, with comparison to
        published work.
  \item Final FYP report, presentation and demonstration.
\end{itemize}
